\section{Future Work}
\label{sec:futureWork}

Several directions can be explored to further improve the approach and evaluate its practical usefulness. First, the current evaluation is based on a single industrial system and a reference dataset of 100 business rules. An important next step is therefore to evaluate the approach on additional industrial systems from different domains. Such evaluations would provide stronger evidence regarding the generalizability of the approach.

Second, the current evaluation measures whether the correct implementation is retained and how much the search space is reduced, but it does not directly measure the effort required by developers to identify the implementation. A controlled developer study would therefore be valuable to assess whether the reduction in candidate methods translates into a reduction in task completion time, cognitive effort, or the number of code elements inspected. Such a study would provide a more direct assessment of the practical benefits of the approach for software maintenance.

Finally, future work could investigate ranking the retained candidate methods according to their likelihood of implementing the business rule. While the current approach deliberately preserves multiple candidates to avoid losing relevant implementations, a ranking mechanism could help developers prioritize their investigation. Other relevant information available from the source code and business rule descriptions could also be explored to improve the ranking of candidate methods. This would move the approach beyond search-space reduction toward more precise and developer-oriented business-rule localization.